% ============================================================================
%  MATEMÁTICAS OPERATIVAS — Sesión 9
%  Semana 9: 5 al 10 de octubre de 2026
%  2.3 Racionalización y sus procesos
%  Duración prevista: 3 horas
% ============================================================================
\documentclass[aspectratio=169,10pt]{beamer}

\usepackage[utf8]{inputenc}
\IfFileExists{spanish.ldf}{\usepackage[spanish,es-noshorthands]{babel}}{}
\usepackage{amsmath,amssymb}
\usepackage{booktabs}
\usepackage{tikz}

\usetheme{Uniminuto}

\uniasignatura{Matemáticas Operativas}
\uniprofesor{Andrés Felipe Puerta Velez}
\unicorreo{andres.puerta-v@uniminuto.edu.co}
\uninrc{CBASBL011}
\unisemana{Sesión 9 — Semana 9 · 5 al 10 de octubre de 2026}

\newcommand{\listatiny}{\setbeamerfont{itemize/enumerate body}{size=\tiny}\tiny}

\begin{document}

\portada

\begin{frame}{Orden del día}
  \begin{enumerate}
    \item Revisión de los ejercicios propuestos de la sesión 8
    \item \azul{Tema 2.3: qué es racionalizar y para qué sirve}
    \item Caso 1: un solo radical en el denominador
    \item Caso 2: índice mayor que 2
    \item Caso 3: binomios y el conjugado
    \item Taller principal, errores frecuentes y cierre
  \end{enumerate}

  \vspace{0.25cm}
  {\footnotesize\color{uniGris} Sesión de 3 horas. Es un tema con
  \destacar{tres recetas} y ninguna teoría difícil. Lo que se evalúa es
  reconocer cuál de las tres aplica, y eso se decide mirando el
  \azul{denominador}.}
\end{frame}

% ===========================================================================
\seccion{2.3 Racionalización y sus procesos}
% ===========================================================================

\begin{frame}{Qué es racionalizar}
  \scriptsize
  \begin{block}{Definición}
    \tiny
    Es transformar una fracción para que el \azul{denominador quede sin
    radicales}, sin cambiar su valor. La herramienta es la misma de la sesión
    3: \destacar{multiplicar arriba y abajo por lo mismo}.
    \[ \frac{a}{b} = \frac{a\cdot k}{b\cdot k} \]
  \end{block}

  \vspace{0.1cm}
  \begin{columns}[T,onlytextwidth]
    \begin{column}{0.48\textwidth}
      \begin{exampleblock}{Por qué se hace}
        \tiny
        Antes de las calculadoras, dividir entre $1{,}732\dots$ a mano era
        penoso; dividir entre $3$ no. Hoy la razón es otra:
        \destacar{la forma racionalizada es la forma canónica} de una
        respuesta, y permite comparar y sumar expresiones que de otro modo
        parecen distintas.
      \end{exampleblock}
    \end{column}
    \begin{column}{0.48\textwidth}
      \begin{alertblock}{Lo que NO cambia}
        \tiny
        \[ \frac{3}{\sqrt{5}} = \frac{3\sqrt{5}}{5} \]
        Las dos valen $1{,}3416\dots$. \azul{Racionalizar no simplifica el
        número, cambia su presentación.} Si el valor cambia, hubo un error.
      \end{alertblock}
    \end{column}
  \end{columns}

  \vspace{0.08cm}
  {\tiny\color{uniGris} \destacar{Verificación gratuita:} calcule con
  decimales la expresión de entrada y la de salida. Si no coinciden, revise.
  Es el control más rápido de todo el tema.}
\end{frame}

\begin{frame}{Caso 1 — un solo radical cuadrático}
  \scriptsize
  \begin{block}{La receta}
    \tiny
    Multiplique arriba y abajo por \destacar{ese mismo radical}. Al cuadrado,
    la raíz desaparece: $\sqrt{a}\cdot\sqrt{a} = a$.
  \end{block}

  \vspace{0.1cm}
  \begin{columns}[T,onlytextwidth]
    \begin{column}{0.48\textwidth}
      \begin{exampleblock}{Ejemplo básico}
        \tiny
        \[ \frac{3}{\sqrt{5}}
           = \frac{3}{\sqrt{5}}\cdot\frac{\sqrt{5}}{\sqrt{5}}
           = \destacar{\frac{3\sqrt{5}}{5}} \]
      \end{exampleblock}
    \end{column}
    \begin{column}{0.48\textwidth}
      \begin{exampleblock}{Con radical arriba y abajo}
        \tiny
        \[ \frac{\sqrt{2}}{\sqrt{7}}
           = \frac{\sqrt{2}\cdot\sqrt{7}}{7}
           = \destacar{\frac{\sqrt{14}}{7}} \]
      \end{exampleblock}
    \end{column}
  \end{columns}

  \vspace{0.12cm}
  \begin{alertblock}{Simplifique el radical antes de racionalizar}
    \tiny
    \[ \frac{6}{\sqrt{12}} = \frac{6}{2\sqrt{3}} = \frac{3}{\sqrt{3}}
       = \frac{3\sqrt{3}}{3} = \destacar{\sqrt{3}} \]
    Si se racionaliza sin simplificar, sale $\tfrac{6\sqrt{12}}{12}$ y
    después hay que simplificar de todos modos. \azul{Menos pasos, menos
    errores.}
  \end{alertblock}
\end{frame}

\begin{frame}{Caso 2 — índice mayor que 2}
  \scriptsize
  \begin{block}{La receta}
    \tiny
    Con $\sqrt[n]{a^{m}}$ en el denominador, hay que multiplicar por
    $\sqrt[n]{a^{\,n-m}}$: así los exponentes suman $n$ y el radical
    desaparece. \destacar{La pregunta es siempre: ¿cuánto le falta al
    exponente para llegar al índice?}
  \end{block}

  \vspace{0.1cm}
  \begin{columns}[T,onlytextwidth]
    \begin{column}{0.48\textwidth}
      \begin{exampleblock}{$\dfrac{5}{\sqrt[3]{2}}$}
        \tiny
        Al exponente $1$ le faltan $2$ para llegar a $3$:
        \[ \frac{5}{\sqrt[3]{2}}\cdot
           \frac{\sqrt[3]{2^{2}}}{\sqrt[3]{2^{2}}}
           = \frac{5\sqrt[3]{4}}{\sqrt[3]{2^{3}}}
           = \destacar{\frac{5\sqrt[3]{4}}{2}} \]
      \end{exampleblock}
    \end{column}
    \begin{column}{0.48\textwidth}
      \begin{exampleblock}{$\dfrac{7}{\sqrt[3]{4}}$}
        \tiny
        $4 = 2^{2}$, y al $2$ le falta $1$:
        \[ \frac{7}{\sqrt[3]{2^{2}}}\cdot
           \frac{\sqrt[3]{2}}{\sqrt[3]{2}}
           = \frac{7\sqrt[3]{2}}{2} \]
        \[ = \destacar{\frac{7\sqrt[3]{2}}{2}} \]
      \end{exampleblock}
    \end{column}
  \end{columns}

  \vspace{0.08cm}
  {\tiny\color{uniGris} \destacar{Error clásico:} multiplicar por
  $\sqrt[3]{2}$ cuando el denominador es $\sqrt[3]{2}$. Eso da
  $\sqrt[3]{4}$, que sigue siendo irracional. \azul{Con índice 3 hacen falta
  tres factores iguales, no dos.}}
\end{frame}

\begin{frame}{Caso 3 — binomios con radicales: el conjugado}
  \scriptsize
  \begin{block}{La herramienta}
    \tiny
    El \azul{conjugado} de $a+b$ es $a-b$. Al multiplicarlos aparece la
    diferencia de cuadrados, que elimina las raíces:
    \[ (a+b)(a-b) = a^{2}-b^{2}
       \qquad\Longrightarrow\qquad
       (\sqrt{7}-\sqrt{3})(\sqrt{7}+\sqrt{3}) = 7-3 = 4 \]
  \end{block}

  \vspace{0.1cm}
  \begin{exampleblock}{Ejemplo desarrollado}
    \tiny
    \[ \frac{6}{\sqrt{7}-\sqrt{3}}
     = \frac{6}{\sqrt{7}-\sqrt{3}}\cdot
       \frac{\sqrt{7}+\sqrt{3}}{\sqrt{7}+\sqrt{3}}
     = \frac{6(\sqrt{7}+\sqrt{3})}{7-3}
     = \frac{6(\sqrt{7}+\sqrt{3})}{4}
     = \destacar{\frac{3(\sqrt{7}+\sqrt{3})}{2}} \]
    \azul{Verificación:} la entrada vale $6\div(2{,}6458-1{,}7321)
    = 6{,}5667$; la salida, $1{,}5\cdot(2{,}6458+1{,}7321) = 6{,}5667$
    \checkmark
  \end{exampleblock}

  \vspace{0.06cm}
  {\tiny\color{uniGris} \destacar{El conjugado solo cambia el signo de en
  medio}, nunca los términos. El conjugado de $2+\sqrt{2}$ es $2-\sqrt{2}$, no
  $-2+\sqrt{2}$ ni $2+\sqrt{2}$.}
\end{frame}

\begin{frame}{Ejercicio 1 — los tres casos}
  \scriptsize
  \begin{alertblock}{En parejas, 12 minutos. Identifique primero el caso.}
    \tiny
    \textbf{1.} $\dfrac{4}{\sqrt{6}}$ \qquad
    \textbf{2.} $\dfrac{10}{\sqrt{20}}$ \qquad
    \textbf{3.} $\dfrac{3}{\sqrt[3]{9}}$ \qquad
    \textbf{4.} $\dfrac{1}{\sqrt{3}+1}$ \qquad
    \textbf{5.} $\dfrac{4}{2+\sqrt{2}}$
  \end{alertblock}

  \vspace{0.05cm}
  \begin{block}{Solución}
    \tiny
    \textbf{1.} $\destacar{\tfrac{4\sqrt{6}}{6}
    = \tfrac{2\sqrt{6}}{3}}$ \qquad
    \textbf{2.} $\tfrac{10}{2\sqrt5} = \tfrac{5}{\sqrt5}
    = \destacar{\sqrt{5}}$ \qquad
    \textbf{3.} $9=3^{2}$, falta un $3$:
    $\tfrac{3\sqrt[3]{3}}{3} = \destacar{\sqrt[3]{3}}$

    \vspace{0.08cm}
    \textbf{4.} $\tfrac{\sqrt3-1}{3-1}
    = \destacar{\tfrac{\sqrt{3}-1}{2}}$ \qquad
    \textbf{5.} $\tfrac{4(2-\sqrt2)}{4-2}
    = \tfrac{4(2-\sqrt2)}{2} = \destacar{4-2\sqrt{2}}$
  \end{block}

  \vspace{0.05cm}
  {\tiny\color{uniGris} Los puntos \destacar{2 y 3} salen ``redondos'' porque
  el radical desaparece del todo. \azul{Eso pasa más de lo que uno espera}, y
  es señal de que el ejercicio se hizo bien.}
\end{frame}

\begin{frame}{Ejercicio 2 — taller principal}
  \begin{alertblock}{En grupos de tres, 20 minutos. Se socializa en el tablero.}
    \scriptsize
    \begin{enumerate}\setlength{\itemsep}{0.0cm}
      \item Racionalice: \ $\dfrac{8}{\sqrt{2}}$
      \item Racionalice: \ $\dfrac{\sqrt{5}}{\sqrt{15}}$
      \item Racionalice: \ $\dfrac{2}{3-\sqrt{5}}$
      \item Racionalice: \ $\dfrac{\sqrt{5}+1}{\sqrt{5}-1}$
      \item Racionalice: \ $\dfrac{1}{\sqrt{x}+\sqrt{y}}$ \ (con $x\neq y$)
      \item Verifique numéricamente el punto 3 con tres decimales.
    \end{enumerate}
  \end{alertblock}
\end{frame}

\begin{frame}{Ejercicio 2 — solución}
  \scriptsize
  \begin{columns}[T,onlytextwidth]
    \begin{column}{0.49\textwidth}
      \begin{block}{Puntos 1 a 3}
        \tiny
        \textbf{1.} $\tfrac{8\sqrt2}{2} = \destacar{4\sqrt{2}}$

        \vspace{0.08cm}
        \textbf{2.} $\tfrac{\sqrt5}{\sqrt{15}}
        = \sqrt{\tfrac{5}{15}} = \tfrac{1}{\sqrt3}
        = \destacar{\tfrac{\sqrt{3}}{3}}$

        \vspace{0.08cm}
        \textbf{3.} $\tfrac{2(3+\sqrt5)}{9-5}
        = \tfrac{2(3+\sqrt5)}{4}
        = \destacar{\tfrac{3+\sqrt{5}}{2}}$
      \end{block}
    \end{column}
    \begin{column}{0.49\textwidth}
      \begin{block}{Puntos 4 a 6}
        \tiny
        \textbf{4.} $\tfrac{(\sqrt5+1)^{2}}{5-1}
        = \tfrac{6+2\sqrt5}{4}
        = \destacar{\tfrac{3+\sqrt{5}}{2}}$

        \vspace{0.08cm}
        \textbf{5.} $\destacar{\tfrac{\sqrt{x}-\sqrt{y}}{x-y}}$

        \vspace{0.08cm}
        \textbf{6.} $\tfrac{2}{3-2{,}236} = \tfrac{2}{0{,}764}
        = 2{,}618$; \ $\tfrac{3+2{,}236}{2} = 2{,}618$ \checkmark
      \end{block}
    \end{column}
  \end{columns}

  \vspace{0.08cm}
  {\tiny\color{uniGris} \destacar{Los puntos 3 y 4 dan lo mismo}, y no es
  casualidad: $\tfrac{\sqrt5+1}{\sqrt5-1}$ y $\tfrac{2}{3-\sqrt5}$ son la
  misma cantidad escrita de dos maneras. \azul{Ese es justamente el argumento
  para racionalizar}: la forma canónica revela igualdades que estaban
  escondidas.}
\end{frame}

\begin{frame}{Errores frecuentes en esta sesión}
  \scriptsize
  \begin{center}
  \begin{tabular}{@{}p{5.2cm}p{6.8cm}@{}}
    \toprule
    \textbf{Error} & \textbf{Lo correcto} \\
    \midrule
    Multiplicar solo el denominador &
      Hay que multiplicar \azul{arriba y abajo}, si no cambia el valor \\
    \addlinespace
    Con $\sqrt[3]{2}$, multiplicar por $\sqrt[3]{2}$ &
      Hacen falta tres factores: multiplique por $\sqrt[3]{2^{2}}$ \\
    \addlinespace
    Conjugado de $2+\sqrt2$ es $-2-\sqrt2$ &
      Es $2-\sqrt{2}$: solo cambia el signo central \\
    \addlinespace
    No simplificar el resultado &
      $\tfrac{6(\sqrt7+\sqrt3)}{4}$ se reduce a $\tfrac{3(\sqrt7+\sqrt3)}{2}$ \\
    \addlinespace
    Racionalizar sin simplificar antes &
      $\tfrac{6}{\sqrt{12}}$: reduzca a $\tfrac{3}{\sqrt3}$ primero \\
    \addlinespace
    $(\sqrt{a}-\sqrt{b})^{2} = a-b$ &
      Es $a - 2\sqrt{ab} + b$; $a-b$ sale del \azul{producto} por el
      conjugado \\
    \bottomrule
  \end{tabular}
  \end{center}
\end{frame}

\begin{frame}{Cierre de la sesión}
  \scriptsize
  \begin{block}{Lo que hay que llevarse}
    \tiny
    \textbf{1.} Racionalizar es cambiar la presentación, no el valor: siempre
    se puede verificar con decimales. \quad
    \textbf{2.} Un radical cuadrático solo: multiplique por él mismo. \quad
    \textbf{3.} Índice $n$: complete el exponente hasta $n$. \quad
    \textbf{4.} Binomio: multiplique por el \destacar{conjugado}, que solo
    cambia el signo del medio. \quad
    \textbf{5.} Simplifique antes y después.
  \end{block}

  \vspace{0.1cm}
  \begin{alertblock}{Para la próxima sesión}
    \tiny
    Resolver los \destacar{20 ejercicios propuestos} y la guía de la unidad en
    Aulas Virtuales.

    \vspace{0.05cm}
    En la sesión 10 vemos \azul{logaritmos y exponenciales}: la operación que
    responde ``¿a qué exponente hay que elevar?''. Es el último tema nuevo de
    la Unidad 2.

    \vspace{0.05cm}
    \destacar{Recuerde: lunes 12 de octubre es festivo} (Día de la Raza).
  \end{alertblock}
\end{frame}

\begin{frame}{Ejercicios propuestos}
  \vspace{-0.1cm}
  \begin{beamercolorbox}[wd=\textwidth,sep=0.10cm]{block body}
    \centering\scriptsize
    Racionalice y simplifique. \quad
    \textbf{1--8:} caso 1. \quad \textbf{9--12:} caso 2. \quad
    \textbf{13--20:} caso 3 y mixtos.
  \end{beamercolorbox}
  \listatiny
  \vspace{0.15cm}
  \begin{columns}[T,onlytextwidth]
    \begin{column}{0.47\textwidth}
      \begin{enumerate}\setlength{\itemsep}{0.16cm}
        \item $\dfrac{5}{\sqrt{3}}$
        \item $\dfrac{7}{\sqrt{7}}$
        \item $\dfrac{12}{\sqrt{6}}$
        \item $\dfrac{9}{\sqrt{18}}$
        \item $\dfrac{\sqrt{3}}{\sqrt{11}}$
        \item $\dfrac{2}{3\sqrt{5}}$
        \item $\dfrac{x}{\sqrt{x}}$
        \item $\dfrac{15}{\sqrt{45}}$
        \item $\dfrac{4}{\sqrt[3]{3}}$
        \item $\dfrac{2}{\sqrt[3]{25}}$
      \end{enumerate}
    \end{column}
    \begin{column}{0.50\textwidth}
      \begin{enumerate}\setlength{\itemsep}{0.16cm}
        \setcounter{enumi}{10}
        \item $\dfrac{6}{\sqrt[4]{8}}$
        \item $\dfrac{1}{\sqrt[5]{4}}$
        \item $\dfrac{1}{\sqrt{5}-2}$
        \item $\dfrac{3}{\sqrt{7}+2}$
        \item $\dfrac{4}{\sqrt{6}-\sqrt{2}}$
        \item $\dfrac{\sqrt{3}}{\sqrt{3}+1}$
        \item $\dfrac{\sqrt{2}+1}{\sqrt{2}-1}$
        \item $\dfrac{5}{2\sqrt{3}-3}$
        \item Escriba el conjugado de $\sqrt{10}-4$
        \item ¿Por qué $(\sqrt a+\sqrt b)(\sqrt a-\sqrt b)$ no tiene
              radicales?
      \end{enumerate}
    \end{column}
  \end{columns}
\end{frame}

\begin{frame}{Respuestas}{Autoevaluación de los ejercicios propuestos}
  \listatiny
  \vspace{0.15cm}
  \begin{columns}[T,onlytextwidth]
    \begin{column}{0.47\textwidth}
      \begin{enumerate}\setlength{\itemsep}{0.16cm}
        \item $\tfrac{5\sqrt{3}}{3}$
        \item $\sqrt{7}$
        \item $2\sqrt{6}$
        \item $\tfrac{3\sqrt{2}}{2}$
        \item $\tfrac{\sqrt{33}}{11}$
        \item $\tfrac{2\sqrt{5}}{15}$
        \item $\sqrt{x}$
        \item $\sqrt{5}$
        \item $\tfrac{4\sqrt[3]{9}}{3}$
        \item $\tfrac{2\sqrt[3]{5}}{5}$
      \end{enumerate}
    \end{column}
    \begin{column}{0.50\textwidth}
      \begin{enumerate}\setlength{\itemsep}{0.16cm}
        \setcounter{enumi}{10}
        \item $3\sqrt[4]{2}$
        \item $\tfrac{\sqrt[5]{8}}{2}$
        \item $\sqrt{5}+2$
        \item $\sqrt{7}-2$
        \item $\sqrt{6}+\sqrt{2}$
        \item $\tfrac{3-\sqrt{3}}{2}$
        \item $3+2\sqrt{2}$
        \item $\tfrac{5(2\sqrt{3}+3)}{3}$
        \item $\sqrt{10}+4$
        \item Porque da $a-b$: la diferencia de cuadrados elimina
              las raíces
      \end{enumerate}
    \end{column}
  \end{columns}
\end{frame}

\begin{frame}{Referencias}
  \footnotesize
  \begin{itemize}
    \item[{[1]}] Ramos del Olmo, Á. M. (2016). \textit{Matemáticas básicas
      para el acceso a la universidad}. Ediciones Pirámide.
      \texttt{https://elibro.net/es/ereader/uniminuto/49134}
    \item[{[2]}] González Pulido, J. W., Chacón Benavides, J. A. \&
      Fonseca Correa, L. Á. (2023). \textit{Matemática básica}
      (1.\textsuperscript{a} ed.). Universidad Mariana.
      \texttt{https://elibro.net/es/ereader/uniminuto/280647}
    \item[{[3]}] Torrijos, M. (2018). \textit{Apuntes de matemática básica}.
      Universidad Católica de Colombia.
      \texttt{https://elibro.net/es/ereader/uniminuto/197048}
    \item[{[4]}] Ramírez, A. \& Rojas, L. (2016). \textit{Matemáticas básicas
      con aplicaciones a la ingeniería}. Ecoe Ediciones.
      \texttt{https://elibro.net/es/ereader/uniminuto/114355}
  \end{itemize}

  \vspace{0.3cm}
  \begin{center}
    {\Large\color{uniAzul}\textbf{¿Preguntas?}}
  \end{center}
\end{frame}

\end{document}
